\chapter{Monads}

\begin{overviewbox}
In Chapter 7 we saw that an adjunction $L \dashv R$ packages a
``construction'' (the left adjoint) together with a ``forgetful''
(the right adjoint). What does this construction \emph{leave behind} on
the source category? The endofunctor $R \circ L : \mathcal{C} \to
\mathcal{C}$, together with the unit $\eta : \mathrm{Id} \Rightarrow R L$
and a ``multiplication'' $\mu = R \varepsilon L : R L R L \Rightarrow R L$
arising from the counit.

This package --- an endofunctor with a unit and a compatible
multiplication --- is called a \term{monad}. Monads are the categorical
formalization of \emph{algebraic structure} carried by a functor. They
appear in algebra (the free-monoid monad, the free-group monad), in
topology (free topological monoids), in logic and computer science (the
list monad, the Maybe monad, the state monad), and in geometry (sheaf
constructions).

The story of this chapter: define monads in their own right, show how
every adjunction gives rise to a monad, study \emph{algebras for a monad}
(the Eilenberg--Moore category), and conclude with the relationship to
the original adjunction. Along the way we develop the Kleisli
construction, the categorical home of effectful computation.
\end{overviewbox}


% ═══════════════════════════════════════════════════════════════════════════
\section{Monads as Algebraic Endofunctors}
\label{sec:monads-defn}
% ═══════════════════════════════════════════════════════════════════════════

\subsection{Informally: Endofunctor + Unit + Multiplication}

A \term{monad} on a category $\mathcal{C}$ consists of three things:
\begin{itemize}
  \item An endofunctor $T : \mathcal{C} \to \mathcal{C}$. Think: ``a way
        of building richer objects from $\mathcal{C}$-objects.''
  \item A unit $\eta : \mathrm{Id}_\mathcal{C} \Rightarrow T$. Think:
        ``the inclusion of $A$ into $T(A)$ as the generators.''
  \item A multiplication $\mu : T \circ T \Rightarrow T$. Think:
        ``compose nested $T$-structures into one.''
\end{itemize}
These satisfy three laws: a unit law (twice) and an associativity law,
to be made precise below.

The classic example: $T(X) = X^*$, the free-monoid (list) functor on
$\mathbf{Set}$. The unit $\eta_X : X \to X^*$ sends $x$ to the singleton
list $[x]$. The multiplication $\mu_X : (X^*)^* \to X^*$ concatenates a
list of lists into a single list.

\subsection{Expanding: Examples That Matter}

\begin{example}{The List Monad on $\mathbf{Set}$}
$T(X) = X^*$ (finite sequences from $X$).
\begin{itemize}
  \item Unit $\eta_X(x) = [x]$.
  \item Multiplication $\mu_X([[a_1, \ldots], [b_1, \ldots], \ldots]) =
        [a_1, \ldots, b_1, \ldots, \ldots]$.
\end{itemize}
The laws say: the unit at the front or back of a list does nothing
($\mu \circ \eta T = \mathrm{id} = \mu \circ T \eta$), and concatenating
nested lists is associative ($\mu \circ T \mu = \mu \circ \mu T$).
\end{example}

\begin{example}{The Maybe Monad}
$T(X) = X + \{*\}$ (a copy of $X$ plus a distinguished ``failure'' value).
\begin{itemize}
  \item Unit $\eta_X(x) = x \in X \subset X + \{*\}$.
  \item Multiplication $\mu_X$ flattens nested options: an outer $*$
        becomes $*$; an outer value that is itself $*$ becomes $*$;
        an outer value that is an inner real value is unwrapped.
\end{itemize}
The Maybe monad is the standard formalism for ``computations that may
fail.''
\end{example}

\begin{example}{The Powerset Monad}
$T(X) = \mathcal{P}(X)$, the powerset.
\begin{itemize}
  \item Unit $\eta_X(x) = \{x\}$.
  \item Multiplication $\mu_X$ takes a set of sets to its union.
\end{itemize}
The Kleisli morphisms (\S 8.4) for this monad are exactly relations:
functions $X \to \mathcal{P}(Y)$ are relations $X \rightharpoonup Y$.
\end{example}

\begin{dialog}
The three examples share a common shape: build something on $X$, ``flatten''
when you build it twice, and have a way to include $X$ into $T(X)$ as the
trivial case.
\end{dialog}

\subsection{Formalizing: The Definition}

A \term{monad on $\mathcal{C}$} is a triple $(T, \eta, \mu)$ where:
\begin{itemize}
  \item $T : \mathcal{C} \to \mathcal{C}$ is a functor.
  \item $\eta : \mathrm{Id}_\mathcal{C} \Rightarrow T$ is a natural
        transformation.
  \item $\mu : T \circ T \Rightarrow T$ is a natural transformation.
\end{itemize}
And satisfying:
\begin{align*}
  \mu \circ \mu T &= \mu \circ T \mu
    && \text{(associativity)} \\
  \mu \circ \eta T &= \mathrm{id}_T
    && \text{(left unit)} \\
  \mu \circ T \eta &= \mathrm{id}_T
    && \text{(right unit).}
\end{align*}
The first equation lives in $T^3 \Rightarrow T$; the others in
$T \Rightarrow T$.

\begin{howtosay}
\begin{itemize}
  \item $T$ --- ``the monad'' (loosely; really it's the whole triple).
  \item $\eta$ --- ``eta,'' the \term{unit} or \term{return}.
  \item $\mu$ --- ``mu,'' the \term{multiplication} or \term{join} (the
        latter in CS).
  \item ``return'' and ``join'' are the names used in functional
        programming for $\eta$ and $\mu$.
  \item The composition $T \circ T$ is sometimes written $T^2$, and the
        $n$-fold composition $T^n$ for any $n$.
\end{itemize}
\end{howtosay}

\subsection{Reorganizing: The Three Laws as Diagrams}

\textbf{Associativity.}
\[
\begin{tikzcd}[column sep=3.5em]
  T^3 \arrow[r, "T \mu"] \arrow[d, "\mu T"'] & T^2 \arrow[d, "\mu"] \\
  T^2 \arrow[r, "\mu"'] & T
\end{tikzcd}
\]

\textbf{Unit laws.}
\[
\begin{tikzcd}[column sep=3em, row sep=2em]
  T \arrow[r, "\eta T"] \arrow[dr, "\mathrm{id}_T"'] & T^2 \arrow[d, "\mu"] & T \arrow[l, "T \eta"'] \arrow[dl, "\mathrm{id}_T"] \\
  & T &
\end{tikzcd}
\]

\subsection{Simplifying: The Monad Schema}

A monad is a \emph{monoid object} in the category of endofunctors of
$\mathcal{C}$ under composition. The unit object of this monoidal
category is $\mathrm{Id}$, the tensor product is functor composition, and
the associativity/unitality of $(T, \eta, \mu)$ are exactly the
associativity/unit laws of a monoid in this monoidal category.

This is why one sometimes hears the slogan ``a monad is a monoid in the
category of endofunctors.'' The slogan is true. Whether it is helpful at
first reading is a matter of taste.

\subsection{Formally: Naturality of $\eta$ and $\mu$}

\formalnote{Naturality is automatic if you write down the right diagrams}
The naturality of $\eta$ says that for $f : A \to B$,
$T(f) \circ \eta_A = \eta_B \circ f$. The naturality of $\mu$ says that
for $f : A \to B$, $T(f) \circ \mu_A = \mu_B \circ T^2(f)$. Both follow
from $\eta, \mu$ being natural transformations between functors of the
appropriate type.


% ═══════════════════════════════════════════════════════════════════════════
\section{Adjunctions Give Monads}
\label{sec:adj-to-monad}
% ═══════════════════════════════════════════════════════════════════════════

\subsection{Informally: $R L$ Is Always a Monad}

Take any adjunction $L \dashv R$ between $\mathcal{C}$ and $\mathcal{D}$.
Set $T = R \circ L$. Then $T$ is an endofunctor of $\mathcal{C}$, and the
adjunction's unit $\eta : \mathrm{Id} \Rightarrow R L = T$ is the unit of
a monad. The multiplication is the whiskering of the counit:
\begin{equation*}
  \mu := R \varepsilon L : R L R L = T^2 \Rightarrow R L = T.
\end{equation*}
The triangle identities of the adjunction translate directly to the
monad laws.

\subsection{Expanding: The Free-Monoid Adjunction $\rightsquigarrow$ The List Monad}

The classical example: $L : \mathbf{Set} \to \mathbf{Mon}$, $R :
\mathbf{Mon} \to \mathbf{Set}$, $L \dashv R$. Then $R L$ on $\mathbf{Set}$
sends a set $X$ to the underlying set of the free monoid on $X$, which is
$X^*$. The induced monad is the list monad.

In general, whenever you see a free/forgetful adjunction, you get a monad
whose algebras (in a sense to be defined below) are the original
structured objects.

\subsection{Formalizing: The Theorem}

\textbf{Theorem.}\quad Let $L \dashv R$ be an adjunction with unit $\eta$
and counit $\varepsilon$. Define $T = R L$ and $\mu = R \varepsilon L$.
Then $(T, \eta, \mu)$ is a monad on $\mathcal{C}$.

\textbf{Proof sketch.} The unit law $\mu \circ \eta T = \mathrm{id}_T$
expands to $R \varepsilon L \circ \eta R L$, which by the triangle
identity for $R$ equals $\mathrm{id}_{R L}$. Similarly for $\mu \circ T
\eta$. Associativity uses the naturality of $\varepsilon$ together with
the fact that $L$ and $R$ are functors. \qed

\subsection{Reorganizing: Every Adjunction Has a Monad}

\begin{itemize}
  \item Free monoid $\dashv$ Forgetful $\rightsquigarrow$ List monad.
  \item Free group $\dashv$ Forgetful $\rightsquigarrow$ Free-group monad.
  \item Discrete $\dashv$ Forgetful $\rightsquigarrow$ Identity monad
        (essentially trivial).
  \item Stone--\v{C}ech $\dashv$ Forgetful (for compact Hausdorff spaces)
        $\rightsquigarrow$ Stone--\v{C}ech monad on $\mathbf{Top}$.
\end{itemize}

\subsection{Simplifying: $T = R L$, $\mu = R \varepsilon L$, $\eta$ as Given}

That is the recipe. Memorize the recipe.

\subsection{Formally: Comonads from Adjunctions}

\formalnote{The dual: comonads}
Just as $R L$ is a monad on $\mathcal{C}$, the composite $L R$ is a
\emph{comonad} on $\mathcal{D}$: a functor with a \term{counit}
$\varepsilon : L R \Rightarrow \mathrm{Id}$ and a \term{comultiplication}
$\delta : L R \Rightarrow L R L R$ satisfying the dual laws. Comonads
formalize ``contexts'' or ``streams'' in computer science.


% ═══════════════════════════════════════════════════════════════════════════
\section{Algebras and the Eilenberg--Moore Category}
\label{sec:em-category}
% ═══════════════════════════════════════════════════════════════════════════

\subsection{Informally: A Monad Is an Algebraic Theory; Algebras Are Its Models}

A monad on $\mathcal{C}$ encodes a kind of algebraic structure. The
\term{algebras for the monad} are the objects of $\mathcal{C}$ carrying
that structure.

Concretely: an algebra for the list monad on $\mathbf{Set}$ is an object
$X$ equipped with a structure map $\alpha : X^* \to X$ ``compatible with
list operations.'' This compatibility makes $\alpha$ a way of
\emph{multiplying} a list of $X$-elements into a single $X$-element ---
i.e., $X$ becomes a monoid. The algebras for the list monad are
\emph{monoids}.

More generally: the algebras for the monad induced by an adjunction
recover the right-hand category, up to equivalence (when nice conditions
hold). This is the \emph{monadicity} story.

\subsection{Expanding: What Is an Algebra?}

\begin{example}{A List-Monad Algebra Is a Monoid}
Given $(X, \alpha : X^* \to X)$ satisfying the laws:
\begin{itemize}
  \item $\alpha \circ \eta_X = \mathrm{id}_X$ (unit): the singleton list
        $[x]$ multiplies to $x$.
  \item $\alpha \circ \mu_X = \alpha \circ T \alpha$ (associativity):
        flattening and then multiplying equals multiplying each list
        and then multiplying the results.
\end{itemize}
From $\alpha$ we recover the monoid operations: the identity is
$\alpha([\,])$ (the empty list), and the binary operation is
$x \cdot y = \alpha([x, y])$. These operations satisfy the monoid
axioms; conversely any monoid yields a list-algebra structure.
\end{example}

\begin{example}{A Powerset-Monad Algebra Is a Complete Sup-Semilattice}
For the powerset monad $T(X) = \mathcal{P}(X)$ on $\mathbf{Set}$, an
algebra structure $\alpha : \mathcal{P}(X) \to X$ assigns to each subset
a supremum, in a way satisfying the laws --- exactly the structure of a
complete sup-semilattice.
\end{example}

\subsection{Formalizing: $T$-Algebras and the EM Category}

For a monad $(T, \eta, \mu)$ on $\mathcal{C}$, a \term{$T$-algebra} is a
pair $(X, \alpha)$ where:
\begin{itemize}
  \item $X$ is an object of $\mathcal{C}$.
  \item $\alpha : T(X) \to X$ is a morphism, the \term{structure map}.
\end{itemize}
Subject to:
\begin{align*}
  \alpha \circ \eta_X &= \mathrm{id}_X
    && \text{(unit)} \\
  \alpha \circ \mu_X &= \alpha \circ T \alpha
    && \text{(associativity).}
\end{align*}

A morphism of $T$-algebras $f : (X, \alpha) \to (Y, \beta)$ is a morphism
$f : X \to Y$ of $\mathcal{C}$ such that $f \circ \alpha = \beta \circ T(f)$.

The category of $T$-algebras and their morphisms is the
\term{Eilenberg--Moore category} of $T$, written $\mathcal{C}^T$ or
$T\text{-}\mathbf{Alg}$.

\subsection{Reorganizing: The EM Adjunction}

The Eilenberg--Moore category comes with a canonical adjunction
$F^T \dashv U^T$:
\begin{itemize}
  \item $U^T : \mathcal{C}^T \to \mathcal{C}$ sends $(X, \alpha)$ to $X$
        and a morphism to its underlying morphism.
  \item $F^T : \mathcal{C} \to \mathcal{C}^T$ sends $X$ to
        $(T(X), \mu_X)$, the \emph{free $T$-algebra} on $X$.
\end{itemize}
The monad induced by $F^T \dashv U^T$ is exactly the original $T$. This
is the \emph{universal} adjunction with this monad: any adjunction giving
rise to $T$ factors uniquely through the EM adjunction.

\subsection{Simplifying: The Algebra Slogan}

Algebras for a monad are the objects on which the monad's structure can
be ``evaluated.''

\subsection{Formally: Monadicity}

\formalnote{Monadicity}
An adjunction $L \dashv R : \mathcal{D} \to \mathcal{C}$ is \term{monadic}
if the canonical comparison functor $K : \mathcal{D} \to \mathcal{C}^T$
(sending $B$ to $(R(B), R \varepsilon_B)$) is an equivalence of
categories. Many ``algebraic'' adjunctions are monadic: monoids over
sets, groups over sets, rings over sets, $R$-modules over abelian groups.
The Beck monadicity theorem gives a precise criterion (which we will not
prove here).


% ═══════════════════════════════════════════════════════════════════════════
\section{The Kleisli Category}
\label{sec:kleisli}
% ═══════════════════════════════════════════════════════════════════════════

\subsection{Informally: Effectful Maps as a New Category}

Many ``computations with effects'' have the type $X \to T(Y)$ rather than
$X \to Y$. A list-valued function $X \to Y^*$. A partial function $X \to
Y_*$. A relation $X \to \mathcal{P}(Y)$. These should be the morphisms
of \emph{some} category. They are: the \term{Kleisli category} of $T$.

\subsection{Expanding: Composition of Effectful Maps}

If $f : X \to T(Y)$ and $g : Y \to T(Z)$, how do we compose them? Not
naively, because the target of $f$ is $T(Y)$, not $Y$. We need to ``lift''
$g$ to $T(Y) \to T^2(Z)$, then flatten via $\mu$:
\begin{equation*}
  g *_K f := \mu_Z \circ T(g) \circ f \;\;:\;\; X \to T(Z).
\end{equation*}
This is \term{Kleisli composition}. The identity at $X$ is the unit $\eta_X
: X \to T(X)$.

\subsection{Formalizing: The Kleisli Category}

For a monad $(T, \eta, \mu)$ on $\mathcal{C}$, the \term{Kleisli category}
$\mathcal{C}_T$ has:
\begin{itemize}
  \item Objects: same as $\mathcal{C}$.
  \item Morphisms $X \to Y$: morphisms $X \to T(Y)$ in $\mathcal{C}$.
  \item Composition: Kleisli composition as above.
  \item Identities: $\eta_X : X \to T(X)$ as the identity at $X$.
\end{itemize}

\textbf{Theorem.}\quad $\mathcal{C}_T$ is a category. The Kleisli
adjunction $F_T \dashv U_T$ between $\mathcal{C}$ and $\mathcal{C}_T$
gives rise to the same monad $T$. Among all adjunctions giving rise to
$T$, the Kleisli adjunction is the \emph{initial} one (every other such
adjunction factors uniquely through it), while the EM adjunction is the
\emph{terminal} one.

\subsection{Reorganizing: EM vs. Kleisli}

Both categories give the same monad $T$. They sit at opposite ends:
\begin{itemize}
  \item Kleisli $\mathcal{C}_T$: ``the smallest category whose adjunction
        yields $T$.''
  \item Eilenberg--Moore $\mathcal{C}^T$: ``the largest such category.''
\end{itemize}
Between them sit all adjunctions giving rise to $T$. The original
adjunction $L \dashv R$ that we started from sits somewhere in this
sandwich.

\subsection{Simplifying: Kleisli for Programmers}

In Haskell-flavoured pseudocode:
\begin{equation*}
  g *_K f = \lambda x.\; \text{``run $f$ on $x$, then run $g$ on the result, flattened.''}
\end{equation*}
This is the \texttt{>>=} (``bind'') operator. The Kleisli category
formalizes ``programs with $T$-effects.''

\subsection{Formally: Kleisli for the List Monad}

\formalnote{Kleisli for lists}
For the list monad, $\mathcal{C}_T = \mathbf{Set}_T$ has objects sets and
morphisms $X \to Y$ given by functions $X \to Y^*$. Kleisli composition
of $f : X \to Y^*$ and $g : Y \to Z^*$ is: apply $f$ to $x \in X$,
yielding a list $[y_1, \ldots, y_n]$; apply $g$ to each $y_i$, yielding
lists $[z_{i1}, \ldots]$; concatenate them all into one list in $Z^*$.
This is the categorical version of nondeterministic computation.


% ═══════════════════════════════════════════════════════════════════════════
\section*{Exercises}
\addcontentsline{toc}{section}{Exercises}
% ═══════════════════════════════════════════════════════════════════════════

\begin{exercisesbox}
\begin{qa}
\qaitem{What is a monad on $\mathcal{C}$?}{A triple $(T, \eta, \mu)$ with $T$ an endofunctor, $\eta : \mathrm{Id} \Rightarrow T$, $\mu : T^2 \Rightarrow T$, satisfying associativity and unit laws.}
\qaitem{What is $\eta$ called?}{The unit (or ``return'').}
\qaitem{What is $\mu$ called?}{The multiplication (or ``join'').}
\qaitem{State the associativity law.}{$\mu \circ \mu T = \mu \circ T \mu$.}
\qaitem{State the two unit laws.}{$\mu \circ \eta T = \mathrm{id}_T = \mu \circ T \eta$.}
\qaitem{Which functor on $\mathbf{Set}$ gives the list monad?}{The free-monoid functor $X \mapsto X^*$.}
\qaitem{What is the unit for the list monad?}{$x \mapsto [x]$.}
\qaitem{What is the multiplication for the list monad?}{Concatenation of lists of lists.}
\qaitem{What is the Maybe monad?}{$T(X) = X + \{*\}$ with $\eta$ the inclusion and $\mu$ the flattening of nested options.}
\qaitem{What is the powerset monad?}{$T(X) = \mathcal{P}(X)$ with $\eta = \{-\}$ and $\mu = \bigcup$.}
\qaitem{How does every adjunction give a monad?}{Set $T = R L$, $\eta$ = adjunction unit, $\mu = R \varepsilon L$.}
\qaitem{Which adjunction gives the list monad?}{Free monoid (left) $\dashv$ Forgetful (right) on $\mathbf{Set} \rightleftarrows \mathbf{Mon}$.}
\qaitem{What is a $T$-algebra?}{A pair $(X, \alpha : T X \to X)$ satisfying $\alpha \circ \eta = \mathrm{id}$ and $\alpha \circ \mu = \alpha \circ T \alpha$.}
\qaitem{What are the algebras for the list monad?}{Monoids.}
\qaitem{What are the algebras for the powerset monad?}{Complete sup-semilattices.}
\qaitem{What is a morphism of $T$-algebras?}{A morphism $f$ of underlying objects such that $f \circ \alpha = \beta \circ T f$.}
\qaitem{What is the Eilenberg--Moore category?}{The category of $T$-algebras and algebra morphisms, denoted $\mathcal{C}^T$.}
\qaitem{What is the free $T$-algebra on $X$?}{$(T(X), \mu_X)$, the algebra given by the monad's own multiplication.}
\qaitem{What adjunction does EM provide?}{$F^T \dashv U^T : \mathcal{C}^T \to \mathcal{C}$, where $U^T$ forgets the structure and $F^T$ is the free algebra.}
\qaitem{What monad does $F^T \dashv U^T$ produce?}{The original $T$.}
\qaitem{What is the Kleisli category $\mathcal{C}_T$?}{Same objects as $\mathcal{C}$; morphisms $X \to Y$ are morphisms $X \to T(Y)$.}
\qaitem{What is Kleisli composition?}{$g *_K f = \mu \circ T(g) \circ f$.}
\qaitem{What is the identity in the Kleisli category?}{$\eta_X : X \to T(X)$.}
\qaitem{What is the role of $\eta$ in the Kleisli category?}{The identity morphism on each object.}
\qaitem{How does EM compare to Kleisli?}{EM is the terminal (largest) adjunction giving $T$; Kleisli is the initial (smallest).}
\qaitem{Is every monad induced by some adjunction?}{Yes, via either EM or Kleisli.}
\qaitem{What is a comonad?}{The dual of a monad: $(W, \varepsilon, \delta)$ with $\varepsilon : W \Rightarrow \mathrm{Id}$, $\delta : W \Rightarrow W^2$.}
\qaitem{Slogan: ``A monad is a monoid in \ldots''?}{``\ldots\ the category of endofunctors with composition as tensor.''}
\qaitem{What does monadicity mean?}{The original adjunction is equivalent to its EM adjunction.}
\qaitem{Why care about Kleisli in computer science?}{Kleisli composition formalizes ``programs with effects''; bind in Haskell is Kleisli composition.}
\qaitem{What is the trivial monad?}{The identity functor with $\eta = \mu = \mathrm{id}$.}
\end{qa}
\end{exercisesbox}


% ═══════════════════════════════════════════════════════════════════════════
\section*{Chapter Summary}
\addcontentsline{toc}{section}{Chapter Summary}
% ═══════════════════════════════════════════════════════════════════════════

\begin{summarybox}
A \textbf{monad} on $\mathcal{C}$ is a triple $(T, \eta, \mu)$: an
endofunctor with a unit and a multiplication satisfying associativity
and unit laws. ``A monoid in the category of endofunctors.''

Every adjunction $L \dashv R$ gives a monad $T = R L$ with multiplication
$\mu = R \varepsilon L$. Conversely, every monad is the monad of some
adjunction --- two canonical choices: \textbf{Eilenberg--Moore} (terminal)
and \textbf{Kleisli} (initial).

A \textbf{$T$-algebra} is a structure map $\alpha : T X \to X$ subject to
the monad laws. The category of $T$-algebras is the Eilenberg--Moore
category $\mathcal{C}^T$. The \textbf{Kleisli category} $\mathcal{C}_T$
has morphisms $X \to T(Y)$ with Kleisli composition.

For the list monad: algebras are monoids, Kleisli morphisms are
list-valued (nondeterministic) functions.
\end{summarybox}


% ═══════════════════════════════════════════════════════════════════════════
\section*{Formal Restatement}
\addcontentsline{toc}{section}{Formal Restatement}
% ═══════════════════════════════════════════════════════════════════════════

\begin{formalrestatement}
\textbf{Monad.}\quad
$(T, \eta, \mu)$ on $\mathcal{C}$ with $T : \mathcal{C} \to \mathcal{C}$,
$\eta : \mathrm{Id} \Rightarrow T$, $\mu : T^2 \Rightarrow T$, satisfying
\begin{align*}
  \mu \circ \mu T = \mu \circ T \mu, \qquad
  \mu \circ \eta T = \mathrm{id}_T = \mu \circ T \eta.
\end{align*}

\medskip
\textbf{Adjunction $\Rightarrow$ monad.}\quad
$L \dashv R$ yields $(R L, \eta, R \varepsilon L)$.

\medskip
\textbf{$T$-algebra.}\quad
$(X, \alpha : T X \to X)$ with $\alpha \circ \eta_X = \mathrm{id}_X$ and
$\alpha \circ \mu_X = \alpha \circ T \alpha$.

\medskip
\textbf{EM category.}\quad
$\mathcal{C}^T$, the category of $T$-algebras and algebra morphisms;
canonically adjoint to $\mathcal{C}$, producing $T$.

\medskip
\textbf{Kleisli category.}\quad
$\mathcal{C}_T$, with $\mathrm{Hom}_{\mathcal{C}_T}(X, Y) = \mathrm{Hom}_\mathcal{C}(X, T Y)$
and Kleisli composition $g *_K f = \mu \circ T(g) \circ f$.

\medskip
\textbf{Slogan.}\quad ``A monad is a monoid in the category of endofunctors.''

\medskip
\noindent\textit{Chapter~9 develops the Yoneda Lemma --- the master
theorem that ties together representable functors, naturality, and
universal properties.}
\end{formalrestatement}
